
%   COMP3310 DV routing tutorial
%   Written by Hugh Fisher, ANU, 2024-
%   Creative Commons CC-BY-NC-SA-4.0 license

%   LaTeX markup for XeTeX with a standard texlive distribution
%   Command to build:
%       xelatex 0N-routing.tex

\input{../common/formatting}

\begin{document}

\TITLE{Distance Vector Routing}

\MINOR{Outline}

In this tutorial you will

\DOT See how distance vector routers self-configure

\DOT See how routers adapt \DSH or not \DSH to changes

\DOT Try sending bad information to routers

\begin{IMPORTANT}
The instructions for this exercise will have even less detail
than usual. Routing is inherently dynamic, unpredictable, and
messy; so the most likely answer to any questions you have is
\DQ{try it and see what happens}.
\end{IMPORTANT}


\SECTION{Setup}

\MINOR{Location}

This exercise is designed for the Computer Science Linux labs.
It is not a stand alone simulator that runs on one machine.
Instead this is designed for a LAN with at least four PCs,
each of which becomes a \SQ{router}.

Ideally there are at least four people. You should be sitting
close together so you can easily see what is happening on
other screens.

If there are less than four, you can log on to more than one
PC in a lab and run the same program on each. On all but one
PC, use \texttt{Ctrl-Alt-F1} before logging on to get
a plain text console, no GUI. (For some reason the keyboard
may stop working on other PCs if you have multiple GUI sessions
open.) After logging out, \texttt{Ctrl-Alt-F7} to restore
the GUI.

The program is written in Python using just the standard
library so could run elsewhere. All of the hosts would need
to be connected to a single Ethernet LAN or a single WiFi access
point, with multicast. Because of this it probably won't work
on the ANU campus WiFi.

You don't need \NAME{WireShark} or \NAME{tcpdump}, as these
\SQ{routers} are not using a standard protocol and will print
out what is happening for you.


\SECTION{Running a router}

Download and unzip the tutorial exercise file. As usual, this is
a CLI program so open a shell and \texttt{cd} into the new directory.
You will find some Python code, a couple of text files, and a
copy of these instructions.

\SECTION{Router links}


\SECTION{Route formation}


\SECTION{Router adaptation}


\SECTION{Misinformation}


\COPYRIGHT

\end{document}
